\section{A common two-dimension Morita reduction}

Five further residual rings admit Morita duals with dimensions
$1$ and $t=2+\sqrt5$, where $t^2=1+4t$. The following arguments
concern arbitrary complex categories. The reductions feed two constructions,
the staged obstruction below, and the final rank-eight endpoint exclusion.

\begin{proposition}\label{prop:four-dihedral}
Every categorification of $R7\_4\_35$, $R7\_4\_2053$,
$R7\_4\_2603$, or $R7\_4\_2665$ is Morita equivalent to a
category with eight invertibles forming $D_8$ and eight noninvertibles
of dimension $t$, on which the invertibles act freely. Its based ring
belongs to the four families described below.
\end{proposition}
\begin{proof}
The source has four invertibles $G$, a simple $X$ of dimension two,
and simples $Y,Z$ of dimension $2t$. There is an involution $h$ fixing
all three on both sides, and
\[
 X^2=\sum_{g\in G}g,\qquad XY=YX=Y+Z,
 \qquad [YY^*:X]=1.
\]
Put $H=\langle h\rangle$. Its pointed associator is trivial because
it fixes $X$, so $A=1+h$ is a connected separable algebra. A fixed
simple supports left and right $A$-module structures, since
$H^2(C_2,\C^\times)=0$. A bimodule supported on a fixed simple $W$
is a projective representation of $H\times H$ on its multiplicity
space. A nontrivial projective class has only two-dimensional
irreducibles. The multiplicity-one channels in $XY$ first force the
classes for $Y,Z$ to be trivial; the $X$ channel in $YY^*$ does the
same for $X$. Thus each of $X,Y,Z$ supports four simple bimodules
with its original underlying object. The two free $H$-orbits in $G$
each support two more. The dual consequently has eight invertibles
and eight simples of dimension $t$. Irrationality and $t^2=1+4t$
imply that $WW^*$ contains exactly one invertible, so the action is free.

Let $K$ be the dual invertible group. The inverse algebra is the sum
of the two invertibles supported on $H$: its multiplicity at a simple
bimodule $V$ is $\dim\operatorname{Hom}_{\mathcal C}(1,V)$.
Call this subgroup $L\leq K$. Its inverse pointed dual is the original
integral subcategory, of rank five and type $1^4,2$.
If $L$ were normal, all four double cosets would have stabilizer $C_2$
and split into two one-dimensional projective representations. That
dual would instead be pointed of rank eight. Hence $L$ is nonnormal,
forcing $K=D_8$ and $L$ to be a reflection subgroup. This does not
trivialize the pointed associator of the dual.

Write $D_8=\langle r,s\mid r^4=s^2=1,\ srs=r^{-1}\rangle$ and
choose a noninvertible $\rho$. Its products have the form
\[
 \rho g=\phi(g)\rho,\qquad \rho^2=a+N\rho,
 \qquad N\in\Z_{\geq0}[D_8],\quad \sum_gN_g=4.
\]
Associativity and duality give
\[
 \phi^2=\operatorname{Ad}_a,\quad\phi(a)=a,\quad
 Ng=\phi(g)N,\quad\phi(N)=N,\quad N^*=a^{-1}N.
\]
The inverse double orbit of $g\rho$ has stabilizer
$L\cap g\phi(L)g^{-1}$. An order-two intersection would give
noninvertibles of dimension $t$ in the source, which has none.
Thus $\phi$ exchanges the two reflection conjugacy classes.
Changing $\rho$ and the generators normalizes
\[
 \phi(r)=r^{-1},\qquad\phi(s)=rs,\qquad a\in\{1,r^2\}.
\]
The coefficient equations give a common coefficient $A$ on rotations
and coefficients $B,B,C,C$ on adjacent reflection pairs, with
$4A+2B+2C=4$. For $a=r^2$, duality also gives $B=C$; a generator
change interchanges $B,C$. There are exactly five possibilities:
\[
\begin{array}{c|c}
a&N\\\hline
1&\sum\text{ rotations},\quad\sum\text{ reflections},\quad2(s+rs)\\
r^2&\sum\text{ rotations},\quad\sum\text{ reflections}.
\end{array}
\]
For $N=2(s+rs)$, the regular Casimir polynomial is
\[
 (x-16)^4(x^2-160x+1280)^2(x^2-96x+768)^4.
\]
Its roots $48\pm16\sqrt6$ force a formal codegree outside
$\mathbb Q(\sqrt5)$. By the published two-dimension Galois reduction
\cite[Lemma~3.1]{Schopieray}, a hypothetical dual can be assumed
pseudo-unitary. The induced-unit formula
\cite[Theorem~2.13]{OstrikRank3} then places every formal codegree
in its simple-dimension field $\mathbb Q(t)$, a contradiction.

The other four families have neutral rank-eight subrings. With $N$
the rotation sum, these are $R8\_1\_15$ for $a=1$ and
$R8\_1\_12$ for $a=r^2$. With $N$ the reflection sum, replacing
$\rho$ by $s\rho$ gives $R8\_1\_18$ in either case. The two
self-dual source cases $35,2053$ require $a=1$; the dual-pair cases
$2603,2665$ require $a=r^2$. The source-bound verifier checks the
full tensors of these statements, including every neutral-subring
product. It does not construct the extensions.
\end{proof}

\begin{proposition}\label{prop:twisted-v4}
Every categorification of $R7\_4\_2591$ is Morita equivalent to a
category with eight invertibles and eight noninvertibles of dimension
$t$, again with a free invertible action.
\end{proposition}
\begin{proof}
Here $G=V_4$ fixes $X,Y,Z$ on both sides, $\dim X=2$,
$\dim Y=\dim Z=2t$, and $Y^*=Z$. The source products include
\[
 X^2=\sum_{g\in G}g,\qquad XY=YX=2Z,\qquad
 XZ=ZX=2Y,\qquad Y^2=2X+4Y+4Z.
\]
The integral subcategory is Tambara--Yamagami. Use its standard
normalization, with nondegenerate symmetric bicharacter $\chi$.
For a $G$-fixed simple $W$, let $\alpha_W,\beta_W\in\mathbb F_2$
be its left and right projective classes and let $B_W$ be its mixed
commutator matrix over $\mathbb F_2$. The full alternating matrix is
\[
 \begin{pmatrix}\alpha_WJ&B_W\\B_W^{\mathsf T}&\beta_WJ\end{pmatrix},
 \qquad J=\begin{pmatrix}0&1\\1&0\end{pmatrix}.
\]
A rank-$2r$ form has projective irreducibles of dimension $2^r$ and
$2^{4-2r}$ inequivalent such irreducibles.

If $W$ occurs $n$ times in $U\otimes V$, its multiplicity space has
diagonal classes $\alpha_U-\alpha_W$, $\beta_V-\beta_W$ and mixed
block $B_W$, up to transpose and inverse. Indeed the left action on
$U$ and right action on $V$ commute before the identifications with
$W$; their residual defect is the mixed commutator of $W$.
The resulting rank is at most $2\log_2n$.
For the TY simple, $\alpha_X=\beta_X=0$ and $B_X=\chi$.
Apply the rank bound to $X$ in $Y^2$. Invertibility of $\chi$ forces
$\alpha_Y=\beta_Y=1$: if either diagonal block vanishes, the rank
is four, whereas both equal to $J$ give rank two because
$\chi^{\mathsf T}J\chi=J$.
Duality gives the same classes for $Z$. Applying the bound to $Z$
in $XY$ gives the matrix
\[
 \begin{pmatrix}J&B_Z\\B_Z^{\mathsf T}&0\end{pmatrix}.
\]
Its rank is $2+\operatorname{rank}(B_Z^{\mathsf T}JB_Z)$, so
$\operatorname{rank}B_Z\leq1$, and likewise for $B_Y$.

Twist the full regular algebra $A_\beta=\bigoplus_{g\in V_4}g$
by the nontrivial element of $H^2(V_4,\C^\times)$. This adds $J$
to both diagonal blocks and leaves the mixed block unchanged.
The $X$ block now has rank two and yields four simple bimodules
with underlying object $2X$, hence dual dimension one.
The $Y,Z$ blocks become
$\left(\begin{smallmatrix}0&B\\B^{\mathsf T}&0\end{smallmatrix}\right)$.
Rank zero would give dual dimension $t/2$, impossible because its
algebraic norm is $-1/4$. Thus both mixed blocks have rank one,
and each yields four simples of dimension $t$.
The pointed support supplies four further invertibles. This accounts
for the whole dual, and the free-action claim again follows from
$t^2=1+4t$.
\end{proof}

The inverse algebra in Proposition~\ref{prop:twisted-v4} is supported
on a $V_4$ subgroup of the dual invertible group. That group can therefore
be $C_4\times C_2$, $C_2^3$, or $D_8$. The staged inverse-algebra argument in
Theorem~\ref{thm:r2591-staged-morita} excludes all these alternatives.
Neither proposition uses a classification restricted to unitary
categories or to trivial pointed associators.

\begin{corollary}\label{cor:all-nine-galois}
Every hypothetical complex categorification of the
rank-seven ring $R7\_4\_2053$ is Galois conjugate to a pseudo-unitary
category.
\end{corollary}
\begin{proof}
For $R7\_4\_2053$ in Proposition~\ref{prop:four-dihedral}, apply
\cite[Lemma~3.1]{Schopieray} to the two-dimension Morita dual.
Categorical global dimension and Frobenius--Perron dimension are Morita
invariant, and the Morita context can be Galois conjugated with the
category. Thus equality of these dimensions transfers back to the
conjugated source. The integral candidate $R7\_3\_601$ is excluded
in Theorem~\ref{thm:r601-center-lattice}.
\end{proof}

The rings $R7\_4\_2665$ and $R7\_4\_35$ are realized in
Theorems~\ref{thm:inversion-orbifold}--\ref{thm:inversion-v4}.
The three quadratic cases are excluded in
Theorem~\ref{thm:quadratic-center-trace}, and $R7\_4\_2591$ is
excluded below. The further two-moment obstruction in
Theorem~\ref{thm:r2603-center-moments} excludes $R7\_4\_2603$.
The final case $R7\_4\_2053$ is excluded by the rank-eight endpoint
in Theorem~\ref{thm:r2053-final}, completing the pinned rank-seven census.
